\section{Correlation independence without a renormalization fixed point}
\label{sec:cpsv16_example_34}

The following example is \cite[Example~3.4]{Cirac2016MPDO_arXiv}. It separates
physical correlation independence from the renormalization fixed-point
condition. It does not assert local orthogonality or zero correlation length in
the stronger basis-of-normal-tensors sense.

\begin{definition}[The tensor of CPSV16 Example 3.4]
    \label{def:cpsv16_example_34_tensor}
    \lean{MPSTensor.cpsvExample34Tensor,
        MPSTensor.cpsvExample34ZeroTensor,
        MPSTensor.cpsvExample34PlusTensor}
    \leanok
    \uses{def:mps_tensor}
    Set $d=D=2$ and
    \begin{align}
        A^0&=\diag\!\left(1,\frac{1}{\sqrt2}\right),
        &
        A^1&=\diag\!\left(0,\frac{1}{\sqrt2}\right).
        \notag
    \end{align}
    The first and second virtual diagonal entries define the bond-dimension-one
    product tensors for $\ket0$ and $\ket+$, respectively.
\end{definition}

\begin{theorem}[Positive-length amplitudes]
    \label{thm:cpsv16_example_34_mpv}
    \lean{MPSTensor.cpsvExample34_mpv}
    \leanok
    \uses{def:cpsv16_example_34_tensor, def:mpv}
    For $N>0$ and $\boldsymbol i=(i_1,\ldots,i_N)\in\{0,1\}^N$,
    \begin{align}
        \Tr(A^{i_1}\cdots A^{i_N})
        &=\mathbf 1_{\{\boldsymbol i=(0,\ldots,0)\}}
          +\left(\frac{1}{\sqrt2}\right)^N.
        \notag
    \end{align}
    Hence the generated MPV is
    $\ket{0}^{\otimes N}+\ket{+}^{\otimes N}$, where
    $\ket{+}=(\ket0+\ket1)/\sqrt2$.
\end{theorem}

\begin{proof}
    \leanok
    \uses{def:mpv}
    Both matrices are diagonal. The first diagonal entry contributes $1$ only
    to the all-zero word, while the second contributes
    $(1/\sqrt2)^N$ to every word.
\end{proof}

\begin{theorem}[Physical correlations are independent of distance]
    \label{thm:cpsv16_example_34_cid}
    \lean{MPSTensor.cpsvExample34_isPhysicalCID}
    \leanok
    \uses{def:cpsv16_example_34_tensor, def:source_physical_cid}
    The tensor in Definition~\ref{def:cpsv16_example_34_tensor} satisfies
    physical correlation independence for arbitrary disjoint physical regions,
    including adjacent regions.
\end{theorem}

\begin{proof}
    \leanok
    Every word matrix is diagonal. Consequently every observable-inserted
    transfer operator and the ordinary transfer operator act entrywise on bond
    matrices, so they commute. The two-region expectation therefore depends on
    the two complementary gap lengths only through their sum.
\end{proof}

\begin{theorem}[Failure of the renormalization fixed-point equation]
    \label{thm:cpsv16_example_34_not_rfp}
    \lean{MPSTensor.cpsvExample34_not_isTransferIdempotent,
        MPSTensor.cpsvExample34_not_hasPhysicalBlockingIsometry}
    \leanok
    \uses{def:cpsv16_example_34_tensor, def:mps_transfer_idempotence,
        def:physical_blocking_isometry}
    The transfer map is not idempotent, and the tensor has no physical blocking
    isometry. Thus it does not satisfy the pure-state renormalization
    fixed-point equation.
\end{theorem}

\begin{proof}
    \leanok
    \uses{thm:physical_blocking_isometry_iff_transfer_idempotence}
    Let $E_{01}=\ketbra0{1}$. The transfer map obeys
    \begin{align}
        \E_A(E_{01})&=\frac{1}{\sqrt2}E_{01},
        &
        \E_A^2(E_{01})&=\frac12E_{01}.
        \notag
    \end{align}
    Hence $\E_A^2\ne\E_A$. The equivalence between transfer idempotence and a
    physical blocking isometry gives the second conclusion.
\end{proof}

\begin{theorem}[Blocked product-state overlap]
    \label{thm:cpsv16_example_34_blocked_overlap}
    \lean{MPSTensor.cpsvExample34_blocked_overlap}
    \leanok
    \uses{def:cpsv16_example_34_tensor, def:mpv_overlap}
    Let $A_0$ and $A_+$ be the bond-dimension-one product tensors obtained
    from the first and second virtual diagonal entries of $A$. Their formal MPV
    overlap at length $n$ is
    \begin{align}
        O_{A_0A_+}(n)
        &=\left(\frac{1}{\sqrt2}\right)^n.
        \notag
    \end{align}
\end{theorem}

\begin{proof}
    \leanok
    \uses{def:mpv_overlap}
    Only the all-zero configuration contributes, with amplitude
    $(1/\sqrt2)^n$.
\end{proof}
